\begin{ouverture}[courbe]{5}{Fonctions et représentations}{Partie 2 · Fonctions}
{Images · Antécédents · Domaines · Modèles}
\objectif{Passer d’une formule à un tableau ou à une courbe.}
\objectif{Déterminer images, antécédents et domaine.}
\objectif{Choisir un modèle continu ou discret.}
\prerequis{Calcul littéral et résolution d’équations.}
\end{ouverture}
\pagestyle{otMain}\essentiel\fontsize{11}{15}\selectfont
\essSection{1}{Une dépendance entre deux grandeurs}
\begin{definition}Une fonction $f$ associe à chaque nombre $x$ de son ensemble de définition $D_f$ un unique nombre $f(x)$, appelé image de $x$. Un antécédent de $y$ est un nombre $x\in D_f$ tel que $f(x)=y$. Un nombre peut avoir plusieurs antécédents, un seul, ou aucun.\end{definition}
\begin{exemple}Pour $f(x)=x^2-1$, $f(3)=8$. Le nombre 8 a deux antécédents : $-3$ et 3. Le nombre $-2$ n’en a aucun, car $x^2\ge0$.\end{exemple}
\begin{savoirfaire}[Déterminer le domaine]
Pour un quotient, le dénominateur doit être non nul. Pour une racine carrée, l’expression sous la racine doit être positive ou nulle. Cumuler les contraintes si plusieurs opérations sont présentes.
\end{savoirfaire}
\begin{exemple}$f(x)=\sqrt{x-2}$ est définie sur $[2;+\infty[$ ; $g(x)=1/(x+3)$ sur $\mathbb R\setminus\{-3\}$.\end{exemple}
\begin{attention}L’image d’un nombre est unique. Chercher les antécédents demande de résoudre une équation et d’en conserver toutes les solutions autorisées.\end{attention}
\clearpage
\essSection{2}{Courbes et changements de registre}
\begin{definition}La courbe $\mathcal C_f$ est l’ensemble des points $(x;y)$ tels que $x\in D_f$ et $y=f(x)$. Ainsi $A(a;b)\in\mathcal C_f$ si et seulement si $a\in D_f$ et $b=f(a)$.\end{definition}
\begin{savoirfaire}[Lire sur une courbe]
Pour lire $f(a)$, partir de $a$ sur l’axe horizontal et rejoindre la courbe verticalement. Pour les antécédents de $b$, chercher toutes les intersections avec la droite horizontale $y=b$, puis lire leurs abscisses.
\end{savoirfaire}
\begin{propriete}Les solutions de $f(x)=g(x)$ sont les abscisses des intersections des deux courbes. Pour $f(x)<g(x)$, repérer les abscisses où la courbe de $f$ est en dessous de celle de $g$.\end{propriete}
\begin{exemple}[Un modèle discret]Le nombre de places occupées après $n$ réservations peut être décrit par $u(n)=2n+5$, avec $n\in\mathbb N$. Les points sont isolés : on ne doit pas attribuer un sens à « 2,4 réservations ».\end{exemple}
\otIntertitre{Avec un tableur}
Pour $f(x)=x^2-1$, placer les abscisses en colonne A ; écrire \code{=A2*A2-1} en B2 et recopier. Le tableau aide à tracer et conjecturer. Il ne prouve pas à lui seul qu’il n’existe aucune autre solution.
\clearpage\rubrique{Découverte}
\begin{blocLivre}{ex-5-1}
\begin{exercice}[1]{Calculer des images}
Pour $f(x)=2x^2-3x+1$, calculer $f(0)$, $f(2)$ et $f(-1)$.
\end{exercice}
\end{blocLivre}
\begin{blocLivre}{ex-5-2}
\begin{exercice}[1]{Lire un tableau}
On donne les valeurs suivantes. Indiquer l’image de $-1$, puis les antécédents de 4 présents dans ce tableau.\begin{center}\renewcommand{\arraystretch}{1.35}\begin{tabular}{c|ccccc}\toprule
$x$ & $-2$ & $-1$ & $0$ & $1$ & $2$\\\midrule
$f(x)$ & $4$ & $1$ & $0$ & $1$ & $4$\\
\bottomrule\end{tabular}\end{center}
\end{exercice}
\end{blocLivre}
\begin{blocLivre}{ex-5-3}
\begin{exercice}[1]{Antécédent d’une fonction affine}
Pour $f(x)=3x-5$, déterminer l’antécédent de 7.
\end{exercice}
\end{blocLivre}
\begin{blocLivre}{ex-5-4}
\begin{exercice}[1]{Appartenir à une courbe}
Pour $f(x)=x^2+2$, les points $A(1;3)$, $B(-2;6)$ et $C(3;10)$ appartiennent-ils à $\mathcal C_f$ ?
\label{fig-ex-5-4}
\figureGeo{fonction-appartenance}
\end{exercice}
\end{blocLivre}
\begin{blocLivre}{ex-5-5}
\begin{exercice}[1]{Le domaine d’un quotient}
Déterminer le domaine de $g(x)=\dfrac{5}{2x-6}$, puis calculer $g(4)$.
\end{exercice}
\end{blocLivre}

\clearpage\rubrique{Entraînement}
\begin{blocLivre}{ex-5-6}
\begin{exercice}[2]{Un domaine avec deux contraintes}
Déterminer le domaine de $h(x)=\dfrac{\sqrt{x+1}}{x-2}$.
\end{exercice}
\end{blocLivre}
\begin{blocLivre}{ex-5-7}
\begin{exercice}[2]{Plusieurs antécédents}
Pour $f(x)=(x-1)^2$, chercher les antécédents de 9, de 0 et de $-1$.
\label{fig-ex-5-7}
\figureGeo{antecedents-parabole}
\end{exercice}
\end{blocLivre}
\begin{blocLivre}{ex-5-8}
\begin{exercice}[2]{Formule et programme de calcul}
Un programme multiplie l’entrée par 4, soustrait 6 puis divise le résultat par 2. Donner sa fonction. Quelle entrée produit 11 ?
\end{exercice}
\end{blocLivre}
\begin{blocLivre}{ex-5-9}
\begin{exercice}[2]{Intersections de deux courbes}
Comparer $f(x)=x^2$ et $g(x)=x+2$ : trouver leurs intersections, puis les abscisses où $f(x)\le g(x)$.
\label{fig-ex-5-9}
\figureGeo{intersection-fonctions}
\end{exercice}
\end{blocLivre}
\begin{blocLivre}{ex-5-10}
\begin{exercice}[2]{Une lecture graphique précise}
La ligne brisée représente une fonction affine sur chaque morceau et passe par $(-2;1)$, $(0;3)$, $(2;-1)$ et $(4;1)$. Lire $f(0)$ et tous les antécédents de 1.
\label{fig-ex-5-10}
\figureGeo{fonction-morceaux}
\end{exercice}
\end{blocLivre}

\clearpage\rubrique{Problèmes}
\begin{blocLivre}{ex-5-11}
\begin{exercice}[3]{Un tarif avec plafond}
Un service facture 3 euros par heure, sans dépasser 15 euros pour la journée. Pour $0\le t\le8$, exprimer le coût $C(t)$ et trouver toutes les durées qui coûtent 12 euros puis 15 euros.
\label{fig-ex-5-11}
\figureGeo{tarif-plafond}
\end{exercice}
\end{blocLivre}
\begin{blocLivre}{ex-5-12}
\begin{exercice}[3]{Une quantité discrète}
Une salle contient 18 chaises avant l’installation de rangées de 6 chaises. On note $u(n)$ le total après $n$ rangées supplémentaires, $n\in\mathbb N$. Donner $u(n)$ ; 75 chaises sont-elles possibles ?
\end{exercice}
\end{blocLivre}
\begin{blocLivre}{ex-5-13}
\begin{exercice}[3]{Deux formules, un même domaine ?}
Comparer les fonctions $f(x)=\dfrac{x^2-1}{x-1}$ et $g(x)=x+1$. Ont-elles le même domaine ? Leurs courbes sont-elles identiques ?
\end{exercice}
\end{blocLivre}
\begin{blocLivre}{ex-5-14}
\begin{exercice}[3]{Concevoir un modèle de volume}
Un réservoir cylindrique de rayon 2 m contient une hauteur d’eau $h$ comprise entre 0 et 3 m. Donner $V(h)$, son domaine et la hauteur correspondant à $6\pi$ m$^3$.
\label{fig-ex-5-14}
\figureGeo{cylindre}
\end{exercice}
\end{blocLivre}
\begin{blocLivre}{ex-5-15}
\begin{exercice}[3]{Une conjecture issue d’un tableau}
Pour $f(x)=x^2-4x+3$, un tableau aux entiers $0,1,2,3,4$ suggère que $f(x)\ge-1$. Démontrer cette inégalité pour tout réel et déterminer le cas d’égalité.
\end{exercice}
\end{blocLivre}

\clearpage
